% GENERATED FILE. Edit canonical lessons, then run npm run generate:books.
% canonical-source-hash: fnv1a64:062adb7dbd8b8870
% canonical-lessons: TE-C344-santapam, TE-C344-maranam, TE-C344-tudisvasa, TE-C344-antimayatra, TE-C344-antyakriyalu

\chapter{Saying Goodbye}
\label{ch:te-saying-goodbye}
\hlchaptermodality{344}

\begin{chapteropening}
\textbf{By the end of this chapter:} \emph{I can say condolence, death, one's last breath, a funeral procession and a funeral.}

Complete the last lesson of chapter 344: I can say condolence, death, one's last breath, a funeral procession and a funeral.
\end{chapteropening}

\section[\texorpdfstring{\te{సంతాపం} (santāpaṁ)}{సంతాపం (santāpaṁ)}]{\te{సంతాపం} (santāpaṁ) — condolence; sympathy}

\label{lesson:TE-C344-santapam}



\begin{quote}
Before the new one: say the Telugu for narrowly, then the Telugu for to get out.
\end{quote}

\subsection*{You'll want to know: \te{సంతాపం}}
\textbf{\te{సంతాపం}} — \emph{santāpaṁ} — \textquotedblleft{}condolence; sympathy\textquotedblright{}.

\begin{quote}
\te{అందరూ} \te{ఆ} \te{కుటుంబానికి} \te{సంతాపం} \te{తెలిపారు}.
\end{quote}

\begin{quote}
\emph{andarū ā kuṭumbāniki santāpaṁ telipāru.}
\end{quote}

\begin{quote}
Everyone offered condolences to that family.
\end{quote}

\te{సంతాపం} \te{తెలుపు}, to express condolence; the one who receives it takes -\te{కి}. It is offered to a family in \te{దుఃఖం}, grief.

\begin{cousinweb}
\textbf{\te{సంతాపం}} (\emph{santāpaṁ}) is borrowed from Sanskrit \emph{saṁtāpa}, \textquotedblleft{}burning heat\textquotedblright{}, and then \textquotedblleft{}affliction, remorse\textquotedblright{} (CDIAL 12885).
\end{cousinweb}

\begin{grammarlens}[title={What you've built}]
One more word for telling what happened.
\end{grammarlens}

\subsection*{Your turn}
\begin{itemize}
\raggedright
  \item \emph{Say it:} \emph{santāpaṁ}
  \item \emph{Say it:} \emph{santāpaṁ}, once more
  \item \emph{Say it:} \emph{andarū ā kuṭumbāniki santāpaṁ telipāru.} — Everyone offered condolences to that family.
  \item \emph{Recall:} say the Telugu for narrowly, then the Telugu for to get out, then say \emph{santāpaṁ} again
\end{itemize}

\begin{tcolorbox}[breakable,colback=teal!4,colframe=teal!35!black,title={Before you move on}]
What does \te{సంతాపం} mean? (\textquotedblleft{}Condolence; sympathy\textquotedblright{}.) Say it once more.
\end{tcolorbox}
\section[\texorpdfstring{\te{మరణం} (maraṇaṁ)}{మరణం (maraṇaṁ)}]{\te{మరణం} (maraṇaṁ) — death}

\label{lesson:TE-C344-maranam}



\begin{quote}
Before the new one: say the Telugu for to get out, then the Telugu for condolence.
\end{quote}

\subsection*{You'll want to know: \te{మరణం}}
\textbf{\te{మరణం}} — \emph{maraṇaṁ} — \textquotedblleft{}death\textquotedblright{}.

\begin{quote}
\te{ఆయన} \te{మరణం} \te{గురించి} \te{పత్రికలో} \te{వార్త} \te{వచ్చింది}.
\end{quote}

\begin{quote}
\emph{āyana maraṇaṁ guriñci patrikalō vārta vaccindi.}
\end{quote}

\begin{quote}
The news of his death appeared in the paper.
\end{quote}

A formal word: in speech people say \te{చనిపోవు}, to die.

\begin{cousinweb}
\textbf{\te{మరణం}} (\emph{maraṇaṁ}) is borrowed from Sanskrit \emph{maraṇa}, \textquotedblleft{}dying, death\textquotedblright{} (CDIAL 9869), from the root \emph{mṛ}, \textquotedblleft{}to die\textquotedblright{}.
\end{cousinweb}

\begin{grammarlens}[title={What you've built}]
One more word for telling what happened.
\end{grammarlens}

\subsection*{Your turn}
\begin{itemize}
\raggedright
  \item \emph{Say it:} \emph{maraṇaṁ}
  \item \emph{Say it:} \emph{maraṇaṁ}, once more
  \item \emph{Say it:} \emph{āyana maraṇaṁ guriñci patrikalō vārta vaccindi.} — The news of his death appeared in the paper.
  \item \emph{Recall:} say the Telugu for to get out, then the Telugu for condolence, then say \emph{maraṇaṁ} again
\end{itemize}

\begin{tcolorbox}[breakable,colback=teal!4,colframe=teal!35!black,title={Before you move on}]
What does \te{మరణం} mean? (\textquotedblleft{}Death\textquotedblright{}.) Say it once more.
\end{tcolorbox}
\section[\texorpdfstring{\te{తుదిశ్వాస} (tudiśvāsa)}{తుదిశ్వాస (tudiśvāsa)}]{\te{తుదిశ్వాస} (tudiśvāsa) — one's last breath}

\label{lesson:TE-C344-tudisvasa}



\begin{quote}
Before the new one: say the Telugu for condolence, then the Telugu for death.
\end{quote}

\subsection*{You'll want to know: \te{తుదిశ్వాస}}
\textbf{\te{తుదిశ్వాస}} — \emph{tudiśvāsa} — \textquotedblleft{}one's last breath\textquotedblright{}.

\begin{quote}
\te{తాతగారు} \te{ఆసుపత్రిలో} \te{తుదిశ్వాస} \te{విడిచారు}.
\end{quote}

\begin{quote}
\emph{tātagāru āsupatrilō tudiśvāsa viḍicāru.}
\end{quote}

\begin{quote}
Grandfather breathed his last in hospital.
\end{quote}

\te{తుదిశ్వాస} \te{విడుచు}, to let go the last breath: a gentle way to say died.

\begin{cousinweb}
\textbf{\te{తుదిశ్వాస}} (\emph{tudiśvāsa}) is a mixed compound: native \emph{tudi}, \textquotedblleft{}the end\textquotedblright{} (DEDR 3314), and Sanskrit \emph{śvāsa}, \textquotedblleft{}breath\textquotedblright{} (CDIAL 12769).
\end{cousinweb}

\begin{grammarlens}[title={What you've built}]
One more word for telling what happened.
\end{grammarlens}

\subsection*{Your turn}
\begin{itemize}
\raggedright
  \item \emph{Say it:} \emph{tudiśvāsa}
  \item \emph{Say it:} \emph{tudiśvāsa}, once more
  \item \emph{Say it:} \emph{tātagāru āsupatrilō tudiśvāsa viḍicāru.} — Grandfather breathed his last in hospital.
  \item \emph{Recall:} say the Telugu for condolence, then the Telugu for death, then say \emph{tudiśvāsa} again
\end{itemize}

\begin{tcolorbox}[breakable,colback=teal!4,colframe=teal!35!black,title={Before you move on}]
What does \te{తుదిశ్వాస} mean? (\textquotedblleft{}One's last breath\textquotedblright{}.) Say it once more.
\end{tcolorbox}
\section[\texorpdfstring{\te{అంతిమయాత్ర} (antimayātra)}{అంతిమయాత్ర (antimayātra)}]{\te{అంతిమయాత్ర} (antimayātra) — a funeral procession}

\label{lesson:TE-C344-antimayatra}



\begin{quote}
Before the new one: say the Telugu for death, then the Telugu for one's last breath.
\end{quote}

\subsection*{You'll want to know: \te{అంతిమయాత్ర}}
\textbf{\te{అంతిమయాత్ర}} — \emph{antimayātra} — \textquotedblleft{}a funeral procession\textquotedblright{}.

\begin{quote}
\te{అంతిమయాత్రలో} \te{వందల} \te{మంది} \te{నడిచారు}.
\end{quote}

\begin{quote}
\emph{antimayātralō vandala mandi naḍicāru.}
\end{quote}

\begin{quote}
Hundreds walked in the funeral procession.
\end{quote}

\te{అంతిమ}, last, and \te{యాత్ర}, a journey: the last journey.

\begin{cousinweb}
\textbf{\te{అంతిమయాత్ర}} (\emph{antimayātra}) is borrowed from Sanskrit: \emph{antima}, \textquotedblleft{}last\textquotedblright{} (Monier-Williams), and \emph{yātrā}, \textquotedblleft{}a journey\textquotedblright{} (CDIAL 10456).
\end{cousinweb}

\begin{grammarlens}[title={What you've built}]
One more word for telling what happened.
\end{grammarlens}

\subsection*{Your turn}
\begin{itemize}
\raggedright
  \item \emph{Say it:} \emph{antimayātra}
  \item \emph{Say it:} \emph{antimayātra}, once more
  \item \emph{Say it:} \emph{antimayātralō vandala mandi naḍicāru.} — Hundreds walked in the funeral procession.
  \item \emph{Recall:} say the Telugu for death, then the Telugu for one's last breath, then say \emph{antimayātra} again
\end{itemize}

\begin{tcolorbox}[breakable,colback=teal!4,colframe=teal!35!black,title={Before you move on}]
What does \te{అంతిమయాత్ర} mean? (\textquotedblleft{}A funeral procession\textquotedblright{}.) Say it once more.
\end{tcolorbox}
\section[\texorpdfstring{\te{అంత్యక్రియలు} (antyakriyalu)}{అంత్యక్రియలు (antyakriyalu)}]{\te{అంత్యక్రియలు} (antyakriyalu) — a funeral; last rites}

\label{lesson:TE-C344-antyakriyalu}



\begin{quote}
Before the new one: say the Telugu for one's last breath, then the Telugu for a funeral procession.
\end{quote}

\subsection*{You'll want to know: \te{అంత్యక్రియలు}}
\textbf{\te{అంత్యక్రియలు}} — \emph{antyakriyalu} — \textquotedblleft{}a funeral; last rites\textquotedblright{}.

\begin{quote}
\te{సాయంత్రం} \te{అంత్యక్రియలు} \te{జరిగాయి}.
\end{quote}

\begin{quote}
\emph{sāyantraṁ antyakriyalu jarigāyi.}
\end{quote}

\begin{quote}
The funeral took place in the evening.
\end{quote}

A plural word: \te{అంత్య}, last, and \te{క్రియలు}, rites or acts.

\begin{cousinweb}
\textbf{\te{అంత్యక్రియలు}} (\emph{antyakriyalu}) is borrowed from Sanskrit: \emph{antya}, \textquotedblleft{}last\textquotedblright{}, and \emph{kriyā}, \textquotedblleft{}work, an act\textquotedblright{} (CDIAL 3587).
\end{cousinweb}

\begin{grammarlens}[title={What you've built}]
One more word for telling what happened.
\end{grammarlens}

\subsection*{Your turn}
\begin{itemize}
\raggedright
  \item \emph{Say it:} \emph{antyakriyalu}
  \item \emph{Say it:} \emph{antyakriyalu}, once more
  \item \emph{Say it:} \emph{sāyantraṁ antyakriyalu jarigāyi.} — The funeral took place in the evening.
  \item \emph{Recall:} say the Telugu for one's last breath, then the Telugu for a funeral procession, then say \emph{antyakriyalu} again
\end{itemize}

\begin{tcolorbox}[breakable,colback=teal!4,colframe=teal!35!black,title={Before you move on}]
What does \te{అంత్యక్రియలు} mean? (\textquotedblleft{}A funeral; last rites\textquotedblright{}.) Say it once more.
\end{tcolorbox}
